\chapter{Nagyobb alkalmazások szervezése}

\noindent\textbf{Tanulási út: Teljes út / skálázódó kódszerkezet.} Akkor olvasd, amikor az egymodulos alkalmazást már nehéz biztonságosan módosítani.\par\medskip

\rustbridgeblockhu{A modulokra bontás, qualified nevek, explicit függőségek és összetartozó kód együtt tartása ugyanazokat a Rustos tervezési ösztönöket használja.}{A Velran application-root source graphot old fel, domain \code{object} namespace-t ad, és a HTTP \code{route} policyt külön kezeli a forrás/modul struktúrától.}{Nem kell újratanulnod a modularitást. Használd a megszokott modulhatárokat; csak a source location, domain namespace és külső HTTP exposure különbségét tanuld meg.}


Ahogy nő az alkalmazás, a change locality fontosabb lesz a fájlszámnál. A \code{mod} a forrásszerkezetet, az \code{object} a domain API névterét szervezi; a HTTP exposure maradjon \code{route}-ban.

\section{Három külön kérdés}

A nagyobb alkalmazás szerkezetét érdemes három külön rétegként kezelni:

\begin{center}
\begin{tabularx}{\textwidth}{@{}l Y@{}}
\toprule
\textbf{Eszköz} & \textbf{Feladata} \\
\midrule
\code{mod} & Melyik fájlban, illetve forrásmodulban található a kód? \\
\code{object} & Melyik üzleti fogalomhoz tartozik a model, query, page vagy action? \\
\code{route} & Milyen HTTP felület és milyen auth/cache/rate-limit policy teszi elérhetővé? \\
\bottomrule
\end{tabularx}
\end{center}

Ez a szétválasztás tudatos. Egy domain object nem kap automatikusan adatbázist, current usert, fájlrendszert vagy hálózati hozzáférést, és egy modul sem jelent HTTP policy-t.

\section{A \code{mod} nem szöveges include}

Egy kisebb, többfájlos projekt már így nézhet ki:

\begin{lstlisting}
myapp/
|-- main.vrn
|-- models.vrn
|-- queries.vrn
|-- components.vrn
|-- pages.vrn
|-- actions.vrn
|-- migrations/
`-- public/
\end{lstlisting}

A belépési pont:

\begin{lstlisting}[language=Velran]
mod models;
mod queries;
mod components;
mod pages;
mod actions;
\end{lstlisting}

A \code{mod} nem C-szerű szövegbeillesztés, és nem önti a deklarációkat globális névtérbe. A compiler az explicit source graphot tölti be, ellenőrzi és egyetlen programmá fordítja. Egy modul saját névteret birtokol: például a \code{queries.vrn} deklarációira más modulból \code{queries::name} alakban kell hivatkozni. A modulútvonal application-root relatív; nested forrásfájlt explicit, például \code{mod article::admin;} deklarációval kell felvenni. Nincs automatikus almappa-discovery vagy relatív \code{../} modulútvonal.

\section{Mikor válts feature-alapú szerkezetre?}

A \code{models.vrn}, \code{queries.vrn}, \code{pages.vrn} típus szerinti bontás jó első lépés, de nagyobb alkalmazásnál könnyen azt eredményezi, hogy egyetlen feature módosításakor öt-hat fájlt kell párhuzamosan megnyitni. Ilyenkor jobb üzleti terület szerint bontani:

\begin{lstlisting}
main.vrn
article.vrn
article/
|-- admin.vrn
`-- components.vrn
account.vrn
account/
`-- profile.vrn
invoice.vrn
invoice/
`-- admin.vrn
routes.vrn
migrations/
public/
\end{lstlisting}

A \code{main.vrn} ilyenkor inkább alkalmazás-kompozíció:

\begin{lstlisting}[language=Velran]
mod article;
mod article::admin;
mod article::components;
mod account;
mod account::profile;
mod invoice;
mod invoice::admin;
mod routes;
\end{lstlisting}

A cél a change locality: egy üzleti változtatás kódja maradjon többnyire egy feature könyvtáron belül.

\section{Domain object: üzleti névtér, nem OOP-osztály}

A globális nevek számát az \code{object} csökkenti. Például:

\needspace{30\baselineskip}
\begin{lstlisting}[language=Velran]
object Article {
    model {
        id: i64
        slug: Slug
        title: String
        body: String
        authorUsername: String
    }

    #[query]
fn bySlug(db: Db, slug: Slug)
        -> Result<Article, DbError> sql {
        SELECT id, slug, title, body, authorUsername
        FROM articles
        WHERE slug = :slug
    }

    #[page]
fn show(ctx: PageContext, db: Db, slug: Slug)
        -> Result<Html, PageError> {
        let article = Article.bySlug(db, slug)?;
        return Ok(html {
            <article>
                <h1>{{ article.title }}</h1>
                @markdown(article.body)
            </article>
        });
    }
}
\end{lstlisting}

A használat így beszédes:

\begin{lstlisting}[language=Velran]
let article = Article.bySlug(db, slug)?;
\end{lstlisting}

Később ugyanilyen természetesek a rövid domainnevek: \code{Invoice.approve}, \code{Invoice.cancel}, illetve \code{UserAccount.disable}.

\subsection*{Mit tartalmazhat egy object?}

A framework domain object egy model köré rendezheti a hozzá tartozó műveleteket. Az alábbi blokk csak szerkezeti váz, nem önmagában fordítható Velran forrás:

\begin{lstlisting}
object Article {
    model { ... }
    #[query]
fn bySlug(...) ...
    #[page]
fn show(...) ...
    #[action]
fn publish(...) ...
    #[component]
fn Card(...) ...
    #[layout]
fn EditorLayout(...) ...
}
\end{lstlisting}

Pontosan egy model blokk tartozik egy framework domain \code{object}-hez. Ez a konstrukció framework-owned model/query/page/action neveket rendez egybe; külön fogalom a normál pure/domain értékekhez használt Rust-szerű \code{struct} + inherent \code{impl} surface-től. A framework object továbbra sem jelent öröklést, rejtett capabilityt, virtual dispatchot vagy reflectiont. Rust-szerű értékeknél támogatott az associated konstruktor (például \code{Type::new}) és az immutable \code{&self} metódus; az általános mutable/owned receiver továbbra is szándékosan fail-closed.


\section{Package-határok újrahasznosítható alkalmazáskódhoz}

Ha kódot több alkalmazás között kell újrahasznosítani vagy stabil library-határ mögé kell zárni, ne növeld korlátlanul az application modulgráfot: emeld lokális package-be. Minden package saját \code{velran.toml}-lal, alapból privát API-val és korlátozott lokális dependency-gráffal rendelkezik. A közvetlen dependency explicit path dependency; a tranzitív dependency enkapszulált marad, nem válik ambient application névtérré.

A library surface tervezéséhez \code{pub struct}, \code{pub fn}, publikus inherent metódus és a szűk package-root modul-re-export használható. A framework authority maradjon az alkalmazásban: egy reusable library \code{pub} hatására sem kaphat HTTP route-ot, actiont, adatbázis-/hálózati authorityt vagy más webes capabilityt. Így a package-dekompozíció architekturális határ, nem capability-eszkalációs mechanizmus.

\section{A dependency továbbra is explicit}

A domain namespace nem rejt el capabilityket. Ez jó:

\needspace{11\baselineskip}
\begin{lstlisting}[language=Velran]
#[query]
fn byId(db: Db, id: i64)
    -> Result<Article, DbError> sql {
    SELECT id, slug, title, body, authorUsername
    FROM articles
    WHERE id = :id
}
\end{lstlisting}

Az action ugyanilyen explicit módon kapja meg a szükséges capabilityket. A következő blokk szándékosan rövidített szerkezeti példa; az állapotváltás konkrét typed queryjét a CRUD és üzleti-audit fejezetek mutatják be:

\needspace{12\baselineskip}
\begin{lstlisting}
#[action]
fn publish(ctx: ActionContext, db: Db, id: i64)
    -> Result<Redirect, PageError> {
    let article = Article.byId(db, id)?;
    authorize article owner authorUsername or role Publisher;
    // ... allapotvaltas tranzakcioban ...
    return Ok(redirect(adminArticles()));
}
\end{lstlisting}

A handler aláírásából továbbra is látszik, hogy adatbázishoz fér hozzá. Ugyanez igaz AppFs-re, authra és named hálózati egressre is: a domain object nem válhat rejtett service locatorrá.

\section{Miért maradjon a route top-level?}

A route tudatosan nem olvad bele teljesen az objectbe:

\begin{lstlisting}[language=Velran]
route articleShow GET "/cikk/:slug<Slug>"
    public cache public ttl 60
    => article::Article.show;

route articlePublish POST "/admin/cikk/:id<i64>/publish"
    auth role Publisher
    invalidate cache articleShow
    => article::Article.publish;
\end{lstlisting}

Így egy route-registry vagy egy code review során egy helyen látszik:

\begin{itemize}
  \item az URL és a HTTP method;
  \item az authentication/authorization követelmény;
  \item a cache policy és invalidáció;
  \item a rate limit vagy resource profile;
  \item a logokban és auditban megjelenő globális route-név.
\end{itemize}

A domain member neve lehet rövid és helyi jelentésű, például \code{Article.publish} a saját modulján belül; külön \code{routes.vrn} fájlból \code{article::Article.publish} szükséges. A route neve viszont maradjon globálisan beszédes, például \code{articlePublish}.

\section{Ajánlott névkonvenció}

\begin{center}
\begin{tabularx}{\textwidth}{@{}l l Y@{}}
\toprule
\textbf{Elem} & \textbf{Forma} & \textbf{Példa} \\
\midrule
object & PascalCase üzleti főnév & \code{Article}, \code{Invoice}, \code{UserAccount} \\
member & camelCase, domainen belül rövid & \code{bySlug}, \code{publish}, \code{approve} \\
route & globálisan auditálható név & \code{articleShow}, \code{invoiceApprove} \\
module & feature vagy technikai egység & \code{article}, \code{account}, \code{routes} \\
\bottomrule
\end{tabularx}
\end{center}

\section{Egy CMS lehetséges szerkezete}

Egy közepes CMS-nél jó kompromisszum lehet:

\needspace{18\baselineskip}
\begin{lstlisting}
main.vrn
routes.vrn
article.vrn            # Article object + alap queryk
article/
|-- admin.vrn          # szerkesztes/publikalas
`-- components.vrn     # Card, listak
media.vrn
media/
`-- admin.vrn
account.vrn
migrations/
public/
`-- app.css
\end{lstlisting}

A publikus és admin HTTP policy továbbra is a route-rétegben marad. Az article feature viszont együtt tartja a cikk üzleti modelljét és műveleteit.

\section{Mit ne csinálj?}

\subsection*{Ne készíts minden deklarációból külön fájlt}

A túl apró modulok több navigációt okoznak, mint amennyi szerkezetet adnak. Egy ötven soros feature-nek nem kell nyolc fájl.

\subsection*{Ne építs rejtett ``service'' réteget}

Ha egy művelet DB-t, AppFs-t vagy auth contextet használ, az maradjon látható a typed aláírásban. A Velran egyik értéke éppen az, hogy a capability-határ explicit.

\subsection*{Ne duplikáld a route policy-t a handlerben}

A route-level auth, cache, rate limit és resource policy deklaratív szerződés. A domain-specifikus objektumszintű authorization természetesen a handlerben maradhat, mert az üzleti rekordtól függ.

\subsection*{Ne használj szabad Stringet üzleti állapotnak}

Ha egy domain \code{Draft}, \code{Review}, \code{Published} vagy \code{Cancelled} állapotokból áll, használd az enumot. A projektstruktúra akkor segít igazán, ha a domain contract is typed.

\section{Refaktorálási út: egy fájlból nagyobb alkalmazás}

Nem kell a projekt első napján végleges könyvtárhierarchiát tervezni. Jó evolúciós út:

\begin{enumerate}
  \item Kezdj egy \code{main.vrn} fájllal, amíg az alkalmazás kicsi.
  \item Válaszd külön a model/query/page/action csoportokat, amikor a fájl már nehezen áttekinthető.
  \item Ha egy feature több technikai fájlon szétszóródik, válts feature-alapú modulokra.
  \item Ha sok globális üzleti név keletkezik, rendezd őket \code{object} namespace-be.
\end{enumerate}

\section{Ellenőrzött companion példák}
Modulstruktúrához használd az \path{examples/module-namespaces/main.vrn} példát, általános többfájlos kiindulóponthoz az \path{examples/starter-project/main.vrn} fájlt, a könyv folyamatos mintaprojektjéhez pedig az \path{examples/secure-notes/main.vrn} példát. Mindhárom compilerrel ellenőrzött, ezért jobb copy/paste forrás a rövidített architekturális vázlatoknál.

\section{Gyakorlat: refaktorálj változási lokalitás alapján}
\paragraph{Gyakorlási mód: önálló.} Használd az előszóban meghatározott önálló módszert.

Válassz egy feature-t, és azonosíts minden fájlt, amely megváltozik a feature módosításakor. Ha egy kis üzleti változás sok, egymástól távoli technikai réteg szerkesztését kívánja, fontold meg a feature vagy domain objektum körüli csoportosítást. Csak akkor refaktorálj, amikor a dependency-irány már világos.

\section{Tipikus hiba: csak a fájlszám miatt modularizálunk}
A hasznos modul koherens contractot birtokol és csökkenti az egyszerre fejben tartandó fogalmakat; kerüld a monolit fájlokat és a mesterséges egyfüggvényes modulokat.

\section{Folyamatos mintaprojekt: Secure Notes}
Mozgasd át a Secure Notes alkalmazást a kezdeti egyfájlos alakból olyan struktúrába, ahol a jegyzethez tartozó modellek, query-k, page/action elemek és policy-k könnyen megtalálhatók. A külső route-ok és security-viselkedés közben ne változzon.
